% !TEX program = lualatex
% Generated by scripts/build_module_practice.py - do not edit by hand.
\DocumentMetadata{
  pdfstandard=UA-2,
  pdfversion=2.0,
  lang=en-US,
  tagging=on
}
\documentclass[11pt]{article}
\usepackage[letterpaper, margin=0.8in, top=0.6in, bottom=0.6in]{geometry}
\usepackage{fontspec}
\setmainfont{Fira Sans}
\usepackage{amsmath}
\usepackage{amssymb}
\usepackage{booktabs}
\usepackage{longtable}
\usepackage{array}
\usepackage{calc}
\usepackage{etoolbox}
\usepackage{xcolor}
\definecolor{accent}{HTML}{BF5700}
\usepackage{titlesec}
\titleformat{\section}{\large\bfseries\color{accent}}{}{0pt}{}
\titlespacing*{\section}{0pt}{14pt}{4pt}
\titleformat{\subsection}{\normalsize\bfseries}{}{0pt}{}
\titlespacing*{\subsection}{0pt}{10pt}{2pt}
\usepackage{hyperref}
\hypersetup{pdftitle={ECON 201 Solutions: The Economist's Toolkit}, pdfauthor={Div Bhagia}, hidelinks}
\pagestyle{plain}
\setlength{\parindent}{0pt}
\setlength{\parskip}{4pt}
\begin{document}
{\LARGE\bfseries\color{accent} Solutions: The Economist's Toolkit}\hfill{\large ECON 201}
\vspace{4pt}


\section{Lecture 3: Making economic decisions: opportunity cost, rents, and incentives}\label{lecture-3-making-economic-decisions-opportunity-cost-rents-and-incentives}

\subsection{1. Adopting a new technology}\label{adopting-a-new-technology}

\textbf{(a)}

\begin{longtable}[]{@{}lrr@{}}
\toprule\noalign{}
& Old press & New press \\
\midrule\noalign{}
\endhead
\bottomrule\noalign{}
\endlastfoot
Benefit & \$150,000 & \$150,000 \\
Direct cost & \$120,000 & \$105,000 \\
\textbf{Net benefit} & \textbf{\$30,000} & \textbf{\$45,000} \\
\end{longtable}

Net benefit is benefit minus direct cost.

\begin{itemize}
\tightlist
\item
  Old press: \$150,000 − \$120,000 = \textbf{\$30,000}
\item
  New press: \$150,000 − \$105,000 = \textbf{\$45,000}
\end{itemize}

\textbf{(b)} The opportunity cost of an option is what Lena gives up by choosing it, which is the net benefit of the next best alternative. She can run only one press, so each option's opportunity cost is the other's net benefit.

\begin{itemize}
\tightlist
\item
  Keeping the old press means giving up the new press, so it costs \textbf{\$45,000}.
\item
  Leasing the new press means giving up the old one, so it costs \textbf{\$30,000}.
\end{itemize}

\textbf{(c)} Economic cost is direct cost plus opportunity cost.

\begin{itemize}
\tightlist
\item
  Old press: \$120,000 of direct cost plus \$45,000 of opportunity cost, so \textbf{\$165,000}.
\item
  New press: \$105,000 of direct cost plus \$30,000 of opportunity cost, so \textbf{\$135,000}.
\end{itemize}

\textbf{(d)} Compare the net benefits of the two options and take the higher: \$45,000 with the new press against \$30,000 with the old one. \textbf{She leases the new press.}

Economic rent is the net benefit of the option chosen minus its opportunity cost.

\begin{itemize}
\tightlist
\item
  New press: \$45,000 − \$30,000 = \textbf{\$15,000 a year}
\end{itemize}

That is her innovation rent, earned by adopting before her rivals.

\textbf{(e)} Redo the net benefits with revenue at \$120,000.

\begin{itemize}
\tightlist
\item
  Old press: \$120,000 − \$120,000 = \textbf{\$0}
\item
  New press: \$120,000 − \$105,000 = \textbf{\$15,000}
\end{itemize}

\textbf{Leasing is still the right choice}, and the rent is still \$15,000 (\$15,000 − \$0). What has fallen is the level of profit: she netted \$45,000 a year while she was the only shop with the new press, and \$15,000 once everyone has it, because competition passed the cost saving on to customers. The reward for moving first was temporary, so firms keep looking for the next technology.

\subsection{2. Is the shop really profitable?}\label{is-the-shop-really-profitable}

\textbf{(a)} The job is the outside option for Lena's own time, so it is the opportunity cost of running the shop, and the shop's net benefit has to beat it.

\begin{itemize}
\tightlist
\item
  Shop with the old press: \$30,000 − \$40,000 = \textbf{−\$10,000 a year}
\end{itemize}

\textbf{No.} The books showed a profit, but in economic terms the shop was losing \$10,000 a year.

\textbf{(b)} Same comparison, with the new press.

\begin{itemize}
\tightlist
\item
  Shop with the new press: \$45,000 − \$40,000 = \textbf{+\$5,000 a year}
\end{itemize}

\textbf{Yes}, and that \$5,000 is her economic rent from running the shop rather than taking the job.

\textbf{(c)} The books count the presses' costs and not Lena's time. Economics counts both: her time has an opportunity cost, the \$40,000 job.

\textbf{(d)} \textbf{Yes, if being her own boss is worth more than \$10,000 a year to her.} The comparison so far counts only money, and running your own shop can be worth something in itself: independence, pride, doing work you chose. The lecture's first move applies here too: put a number on it, the most Lena would pay to be her own boss rather than an employee. If that number is more than the \$10,000 gap, keeping the shop passes the decision rule after all.

\subsection{3. Multiple choice}\label{multiple-choice}

\medskip

\textbf{1.} \textbf{(c).} The plane, crew, and fuel are paid for whether or not the seat is filled, so almost nothing is given up by letting one more passenger on.

\medskip

\textbf{2.} \textbf{(b).} Going to Band A means giving up the evening with Band B, and the net benefit of that alternative is \$50 − \$40 = \textbf{\$10}.

\section{Lecture 4: Why do markets matter? Specialization and comparative advantage}\label{lecture-4-why-do-markets-matter-specialization-and-comparative-advantage}

\subsection{1. Pablo and Lin strike different deals}\label{pablo-and-lin-strike-different-deals}

\textbf{(a)} Lin receives 1.5 × 2 = 3 loaves.

\begin{itemize}
\tightlist
\item
  Lin: \textbf{3 loaves and 2 pounds}, against 1.5 loaves and 1 pound on her own. Better off in both goods.
\item
  Pablo: 12 − 3 = \textbf{9 loaves and 2 pounds}, against 9 loaves and 1.5 pounds on his own. The same bread and half a pound more cheese, so better off.
\end{itemize}

\textbf{(b)} \textbf{Lin's.} Both prices sit inside the mutually beneficial range, between 0.5 and 2 loaves per pound, but 1.5 is closer to Pablo's opportunity cost of 2. The nearer the price is to the buyer's own cost of producing the good, the more of the gain goes to the seller, and Lin is the seller here.

\textbf{(c)} \textbf{No trade happens.} Making a pound of cheese costs Lin half a loaf of bread, so selling it for a quarter of a loaf leaves her worse off than baking her own bread. A price below the seller's opportunity cost kills the deal.

\subsection{2. Two food trucks}\label{two-food-trucks}

\textbf{(a)}

\begin{longtable}[]{@{}lrr@{}}
\toprule\noalign{}
Opportunity cost of\ldots{} & 1 burrito & 1 taco \\
\midrule\noalign{}
\endhead
\bottomrule\noalign{}
\endlastfoot
Luis & 40 / 8 = 5 tacos & 8 / 40 = 0.2 burritos \\
Rosa & 12 / 6 = 2 tacos & 6 / 12 = 0.5 burritos \\
\end{longtable}

\textbf{(b)} \textbf{Luis has the absolute advantage in both} (40 tacos against 12, and 8 burritos against 6). But a burrito costs Rosa only 2 tacos against Luis's 5, so \textbf{Rosa has the comparative advantage in burritos}, and a taco costs Luis 0.2 burritos against Rosa's 0.5, so \textbf{Luis has the comparative advantage in tacos}.

\textbf{(c)}

\begin{longtable}[]{@{}lrr@{}}
\toprule\noalign{}
& Tacos & Burritos \\
\midrule\noalign{}
\endhead
\bottomrule\noalign{}
\endlastfoot
Luis & 0.75 × 40 = 30 & 0.25 × 8 = 2 \\
Rosa & 0.67 × 12 = 8 & 0.33 × 6 = 2 \\
\textbf{Total} & \textbf{38} & \textbf{4} \\
\end{longtable}

\textbf{(d)} Luis makes only tacos, 40 of them, and Rosa makes only burritos, 6 of them. Compared with 38 tacos and 4 burritos under self-sufficiency, specialization yields \textbf{2 more tacos and 2 more burritos} from the same two people and the same hours.

\textbf{(e)} The price is 9 / 3 = 3 tacos per burrito.

\begin{itemize}
\tightlist
\item
  Luis: 40 − 9 = \textbf{31 tacos} and the \textbf{3 burritos} he bought, against 30 and 2 on his own. Better off.
\item
  Rosa: the \textbf{9 tacos} she received and 6 − 3 = \textbf{3 burritos}, against 8 and 2 on her own. Better off.
\end{itemize}

Everyone ends up with more of both products.

\textbf{(f)} Rosa sells only above her own cost of a burrito, 2 tacos; Luis buys only below his, 5 tacos. Any price \textbf{between 2 and 5 tacos per burrito} leaves both better off, which is why the deal at 3 worked.

\subsection{3. Ana and Ben at the campus bakery}\label{ana-and-ben-at-the-campus-bakery}

\textbf{(a)} \textbf{Ana in frosting} (12 cupcakes against 4) and \textbf{Ben in coffee} (4 pots against 3).

\textbf{(b)} Work out each person's opportunity cost, what they give up. The hour Ana spends brewing 3 pots could have frosted 12 cupcakes instead, so each pot costs her 12 / 3 = 4 cupcakes. The same hour costs Ben only 4 / 4 = 1 cupcake per pot, so \textbf{Ben has the comparative advantage in coffee}. Flipping it around, a cupcake costs Ana 3 / 12 = 0.25 pots and Ben 1 pot, so \textbf{Ana has the comparative advantage in frosting}. Here comparative advantage lines up with absolute advantage, which happens whenever each person is better at one task.

\textbf{(c)} Ben only gains at a price above his own cost of making a pot, 1 cupcake. Ana only gains at a price below what a pot costs her to brew herself, 4 cupcakes. So both gain at any price \textbf{between 1 and 4 cupcakes per pot}. Where the price lands inside that range decides who captures more of the gains: near 1 favors Ana, near 4 favors Ben.

\subsection{4. Multiple choice}\label{multiple-choice}

\medskip

\textbf{1.} \textbf{(b).} The \$45 is a sunk cost: it has been incurred and cannot be recovered, so it is the same under both choices and should not sway the decision. Options (c) and (d) are incorrect because they ignore the opportunity cost of her evening, the night of rest she values.

\medskip

\textbf{2.} \textbf{(b).} The time that makes 5 necklaces could instead make 20 bracelets, so each necklace costs 20 / 5 = \textbf{4 bracelets}.

\medskip

\textbf{3.} \textbf{(a) and (d).} Jose produces more of both goods, so (a) is true and (b) is false. For comparative advantage, look at where each person's edge is biggest: Jose is twice as good as Alex at oranges, but only 1.5 times as good at melons, so Jose's comparative advantage is in oranges and Alex's is in melons, which makes (d) true. And an orange costs Jose half a melon, while it costs Alex two thirds of one, so Jose's opportunity cost of an orange is lower, not higher: (c) is false.

\medskip

\textbf{4.} \textbf{(b).} Jose has the comparative advantage in oranges, so he is the orange seller. A higher price in melons per orange means every orange he sells brings home more melons, so higher prices shift the gains toward Jose.

\medskip

\textbf{5.} \textbf{(c).} Comparative advantage is about relative costs, and relative costs cannot be lower across the board: being relatively cheaper at one good means being relatively more expensive at another. So it still pays to concentrate where the advantage is largest and trade for the rest.

\medskip

\textbf{6.} \textbf{(c).} Since Ravi's opportunity cost of coffee is lower, he is the one selling coffee and Sofia is buying. Ravi only gains at a price above his own cost of 2 kilos of rice, and Sofia only gains below her cost of 5, so both gain at any price strictly between 2 and 5.

\medskip

\textbf{7.} \textbf{(c).} Typing it herself is not free: those hours have an opportunity cost, and for her they are worth far more in the studio or on tour than an assistant's wage.

\medskip

\textbf{8.} \textbf{(b).} Trade patterns follow comparative advantage, not absolute advantage. Even if an American factory could sew a shirt in fewer hours, those hours earn far more in the industries where the US advantage is greatest, so shirts are cheaper to buy from abroad than to make at the cost of that forgone output.

\medskip

\textbf{9.} \textbf{(b) and (c).} Specializing means depending on others, so a disruption to trade can leave you exposed, and a producer of one thing is vulnerable if demand for it falls. Specialization raises total output and both parties can gain, so (a) and (d) are false.

\end{document}
